\scanpage{78}
\uline{Relations avec op.} $\otimes$ \uline{et} $\Gamma^p$.

La donnée d'une cat. polynomiale \inserted{$P$} permet
de définir un foncteur triadditif
\[
P_1^\circ\times P_1^\circ\times P_1\longrightarrow\mathrm{Ab}
\]
par
\[
(X,Y;Z)\longmapsto P_{(1,1)}(X,Y;Z)
\]
(et de $=$ pour plus d'\uline{arguments} à gauche)

Si, pour $X,Y$ fixés, $Z$ variable, ce foncteur
est représentable, l'objet qui le représente est
noté $X\otimes Y$. Il dépend fonctorielt de
$X,Y$ ds $P_1$ (mais non pour $X,Y\deleted{\in\operatorname{Ob}P}$ considérés
comme dans $P$ !). \deleted{\unsure{On}} voit que $X\otimes Y$ existe
ssi $Y\otimes X$ existe, et ils sont can. isom. (\deleted{\unsure{fl|par}}
$s_{X,Y}$), on a $s_{X\otimes Y}\cdot s_{Y\otimes X}=\mathrm{id}_{Y\otimes X}$.

Soit $X,n\in\mathbb N$ fixés, si le foncteur
\[
Y\longmapsto P_n(X,Y)
\]
sur $P_1$ est représentbl, l'objet qui le représente
est noté $\Gamma^n X$ ; il \inserted{en} dépend aussi fonctorielt
de $X$ ds $P_1$. L'objet $\unsure{\Gamma_{\scriptscriptstyle\circ}^0|\Gamma^0}(X)$ existe ssi $\Gamma^0(0)$ existe
il est noté $\mathbb{1}$.

\unsure{Prouvons|Pouvons} que \deleted{\unsure{$\forall$}} si $\mathbb{1}$ existe, alors $\forall X\in\operatorname{Ob}\unsure{P|\Gamma|\mathcal P}$,
$\mathbb{1}\otimes X\simeq X$ i.e.
\[
P_{(1,1)}(\mathbb{1},X;Y)\underset{\inserted{\text{isom bifonctoriel}}}{\simeq}P_1(X,Y)(=\operatorname{Hom}_{P_1}(X,Y))\quad ?
\]
et plus particulièrement $\mathbb{1}\otimes\mathbb{1}\simeq\mathbb{1}$ i.e. $\forall Y$
\[
P_{1,1}(\mathbb{1},\mathbb{1};Y)\simeq P_1(\mathbb{1},Y)\quad ?
\]
% archive sheet number: 77
